% !TEX root = ../main.tex
\section{Introduction}

Linear programming has been the workhorse of strategic forest planning for half a century. Since the Model~I and Model~II formulations of \citet{johnson1977techniques}, LP harvest-scheduling models, such as the USDA Forest Service's FORPLAN system and its successors, have been used to project long-term timber supply, to test harvest-flow policies, and to support trade-off analysis through the shadow prices that an optimal LP solution provides. In nearly all of such applications, the model is solved once, over the full multi-decade planning horizon, as an \emph{open-loop} problem: the entire sequence of future decisions is chosen at time zero, and the resulting schedule of harvest levels and residual forest structure is presented as the plan.

Real-world forest planning likely does not align with the open-loop assumption that today's planner could precommit to every future decision. Real planning is sequential: each period, the planner observes the realized forest state and re-solves the planning problem from that state, under the same objectives and constraints. If the tail of the original plan is not optimal for the subproblem that starts at the realized state, the re-solving planner will deviate from it. The announced plan then fails Bellman's principle of optimality \citep{bellman1957dynamic} and is \emph{dynamically inconsistent}. The projected harvest trajectory, residual growing stock, and shadow prices are not credible statements about what will or should happen. A plan that will not be followed cannot be a sound basis for long-term sustainable policies. Comparative analyses performed on it, including shadow-price trade-off analyses that often motivate such modelling studies, are biased.

This trap is a special case of a general phenomenon in dynamic economics: a plan that is optimal at time zero need not remain optimal at a later time, so a sequence of planners each re-optimizing from the realized state will not follow it \citep{strotz1956myopia,kydland1977rules}. It was demonstrated carefully and thoroughly for forest planning more than three decades ago. In his 1991 PhD thesis at the University of California, Berkeley, P. J. Daugherty proved analytically and demonstrated numerically that LP-based forest-planning models solved as open-loop formulations admit dynamically inconsistent plans, identified the structural conditions under which the inconsistency arises, and characterized its behaviour over a wide range of initial forest conditions and harvest policies \citep{daugherty1991credibility}. The thesis remains, to our knowledge, the most complete treatment of the problem in the forest-planning literature, yet it was never published in the peer-reviewed literature and is held only as an archival hard copy. The phenomenon itself has surfaced in the literature under various names, before and since---\citet{mcquillan1986declining} named the ``declining even-flow effect'' in national forest planning; \citet[pp.~331--332]{gunn2007models}, citing \citet{daugherty1991credibility}, calls it the ``declining non-declining yield''; and \citet{paradis2013risk} documented the systematic drift that arises between planned and implemented harvest under incoherent hierarchical planning, citing \citet{daugherty1991credibility} directly and using ``credibility'' in the thesis's sense. Yet the original analysis remains largely unread, and the field has lacked an open, citable, reproducible reference implementation against which to check it.

This paper formalizes dynamic inconsistency for LP forest plans and provides that reference. We present an open, transparent, and fully reproducible reproduction and refinement of the modelling methods and case study of \citet{daugherty1991credibility}, built on the open-source \texttt{ws3} wood-supply model \citep{paradis2026ws3} from the public-domain Umpqua planning documents (the 1987 Draft EIS and the 1990 FEIS Appendix B yield tables and economics). The contribution is fourfold. First, we give an operational divergence metric and occurrence criterion for dynamic inconsistency, making the phenomenon measurable and comparable across models. Second, we provide an open, test-backed \texttt{ws3} stack---the case-study data, the LP formulations, the sequential-replanning simulator, and the experiment grid---including a new \texttt{ws3} constraint form for the classic FORPLAN harvest-flow policies (non-declining yield and bounded sequential flow), released in \texttt{ws3} v1.1.0a5. Third, we reproduce and refine the thesis's empirical characterization across its full experiment grid (18 initial forest conditions $\times$ 4 discount rates $\times$ 6 harvest-flow policies), cleanly separating the harvest-flow constraint's role against a no-flow control and recovering the discount rate's occurrence-invariance and magnitude-dependence. Fourth---moving beyond reproduction---we run four pre-registered-style extension experiments, each asking the same question of a candidate modification: does it mitigate, or eliminate, the inconsistency? The variants probe the discount-rate \emph{shape} (declining-rate schedules), the flow constraint's \emph{denomination} (net revenue instead of volume) and \emph{shape} (rolling-mean rather than pointwise linkages, under two replanning institutions), and a \emph{cap-calibration} alternative that replaces the flow link with an iteratively tightened max-harvest cap. We emphasize the framing: the underlying science is a faithful reproduction of \citet{daugherty1991credibility}; the novelty is the operational metric, the open and reproducible benchmark, the rigorous driver diagnosis, and the systematic mitigation tests the open stack makes possible.

The remainder of this paper is organized as follows. Section 2 reviews the theory of dynamic inconsistency and the open-loop LP forest-planning problem, including the Model I/II/III formulation distinction. Section 3 describes the case study, the case-study data, the sequential-replanning simulator, the inconsistency metric, and the experiment design. Section 4 presents the reproduced results and the four extension experiments. Section 5 discusses the drivers of the inconsistency, the discount-rate nuance, implications for the literature, remedies, and limitations. Section 6 concludes.
